\section{Roadmap Update after Experiment 13B}
\label{sec:roadmap-post-13b}

Experiment~13B removes the immediate need for another synthetic-network communication mechanism. The raw-threshold LIF rule survives two hidden layers, covers essentially all parameter tensors, scales from 169 to more than 600 parameters, and remains strongly communication-efficient relative to EF-TopK and dense asynchronous SGD. The experiment also shows that the apparent deep-network loss gap must be interpreted carefully: most of the large gap in the regularized test metric was caused by comparing different parameter norms and dimensions. Predictive cross-entropy reveals only a small depth penalty, shared by all optimizers.

Two negative findings should shape the next phase. First, simply inserting Experiment~13A's very small initialization into an extra $\tanh$ layer causes a method-independent vanishing-signal failure. Future neural-network experiments must therefore use standard trainability checks before interpreting communication failures. Second, there is no evidence that layer starvation currently requires layer-wise thresholds, block competition, or microblock event coding. Adding such mechanisms now would increase algorithmic complexity without addressing an observed limitation.

The next scientific gate should therefore move away from increasingly elaborate synthetic MLPs and test whether the event mechanism survives a more realistic data/model setting. A suitable Experiment~14A is a real-image federated classification benchmark with deliberately controlled non-IID client partitions. A small MLP or compact convolutional network should be used first, with full precision and EF-TopK retained as communication references. The primary metrics should be predictive test loss/accuracy versus cumulative bits, event density per layer, parameter coverage, and robustness across label-skew strengths.

Clustering/personalization should remain a later branch unless the real-data experiment exhibits a clear multimodal-client failure. If a single global event-trained model remains competitive under ordinary label skew, adding clustering would solve no demonstrated problem. If instead client groups exhibit incompatible optima or strongly different event-sign/co-firing statistics, a later experiment can ask whether the event stream itself provides a useful low-communication similarity signal for clustered or personalized FL.

The resulting progression is
\begin{equation}
 \boxed{
 13\mathrm{B}\;\text{depth/dimension PASS}
 \;\rightarrow\;
 14\mathrm{A}\;\text{real-data scaling}
 \;\rightarrow\;
 \text{mechanism or clustering only if a concrete failure appears}.
 }
\end{equation}
